\section{Cómo la formación en física moldeó el gusto por la investigación en IA}

La experiencia de Yao Shunyu (姚顺宇) en física no es un detalle incidental de la entrevista, sino un antecedente importante para entender su criterio de investigación. En la licenciatura trabajó en teoría de la materia condensada, y después en física teórica de altas energías, información cuántica y temas relacionados con agujeros negros. Al hablar de los sistemas no hermitianos, subrayó que la investigación de su etapa de licenciatura tiene semejanzas con lo que hace ahora en IA: primero se tiene una comprensión o una idea, después se diseña un experimento numérico para verificarla y, por último, la idea se concreta en un pipeline.

\begin{figure}[H]
\centering
\includegraphics[width=0.88\textwidth]{frame-physics-nonhermitian.jpg}
\caption{Sección sobre la experiencia en física: de la materia condensada y los sistemas no hermitianos al método experimental en IA.\protect\footnotemark}
\end{figure}
\footnotetext{Intervalo de tiempo del fotograma del video: 01:28:25--01:28:55. Este fotograma se usa para marcar la discusión sobre el trasfondo de física.}

\subsection{Explicación intuitiva de los sistemas no hermitianos}

En la mecánica cuántica ordinaria, el hamiltoniano suele ser hermitiano, lo cual garantiza que los valores propios de energía sean reales y que la evolución del sistema conserve la probabilidad. Los sistemas no hermitianos, en cambio, suelen usarse para describir sistemas abiertos, disipación, ganancia, medición o dinámicas efectivas. Para quien lee estas notas de IA no es necesario profundizar en los detalles físicos, pero sí conviene entender el hábito de investigación que esto formó: ante un sistema complejo, primero se construye un modelo verificable y después se usa un experimento numérico para acercarse a su comprensión.

\begin{knowledgebox}{La transferencia de la física a la IA}
La formación en física enfatiza: definir el problema, construir un modelo simplificado, diseñar el experimento, observar las anomalías y corregir la teoría. El entrenamiento de los modelos grandes es similar: se plantea una hipótesis sobre el comportamiento, se diseñan los datos y el pipeline de entrenamiento, se observan mediante experimentos la pérdida, los benchmarks, el comportamiento de los usuarios y los casos de falla, y después se vuelve a la hipótesis. La diferencia es que los sistemas de IA tienen un carácter más ingenieril y dependen más de la colaboración organizacional.
\end{knowledgebox}

\subsection{«Retar lo que no domina»}

En la entrevista, Yao Shunyu dijo varias veces que le gusta retarse con cosas en las que no es bueno. Dejar la física no fue porque la física le pareciera aburrida, sino porque sentía que en esa dirección ya había visto a mucha gente más inteligente que él, y también quería retarse en un campo nuevo. Esta decisión explica la motivación que tuvo después para ir a Anthropic y luego a Gemini: no era buscar un título, sino encontrar un entorno que lo obligara a aprender cosas nuevas.

\begin{warningbox}{No mitificar el cambio de campo}
Pasar de la física a la IA no implica automáticamente una especie de «ventaja aplastante». Yao Shunyu, por el contrario, dijo varias veces que la IA no necesariamente requiere la inteligencia más abstracta, sino que requiere ser confiable, meticuloso y responsable. El valor que aporta el cambio de campo es principalmente el gusto de investigación y los hábitos experimentales, no una superioridad innata.
\end{warningbox}

\subsection{Resumen de la sección}

El trasfondo en física le dio a Yao Shunyu una manera de abordar sistemas complejos: de la idea al experimento, del experimento a la comprensión, y de ahí al pipeline. Al transferirse a la IA, este método se manifiesta en la importancia que da a la definición del problema, la construcción de datos y la verificación experimental.
